\documentclass{article}
\usepackage{amsmath}
\begin{document}
\section*{Preserve the prepared read, not an ambient execution assumption}
A prepared diagnostic is for a named stored result, one physical
source and one receiver. The tested sender and receiver own distinct
native allocations. The scheduling law is that unrelated prepared
work should keep its scope while a different pair of transmitters
is replaced; changed local authority, a changed allocation or allowed
audience, or a newly binding resource constraint must still stop a
send. The test does not equate local scope with a claim of independent
service. It separately asks whether the native broker can execute
an admitted intent at the cut.

\paragraph{The two crossings.}
First compare evidence: the exact send intent descends to both
contexts if the named content, governing law, source and receiving
allocations, and singleton audience agree. The implementation
compares the current native allocations with the retained
allocations, not their current string names. An ambient change to
the membership generation is not a counterexample to that
comparison. Replacing either endpoint by a different native
allocation, revoking the applicable law, or adding a subscriber
does break the scoped comparison. Extra recipients change both
the act and its debit footprint. An old exact receipt remains an
immutable value even when its prepared forwarding no longer
qualifies.

Second compare action: admissibility of the carried send uses
the live charge profiles and broker readiness. The selected native
dispatch path updates the poses of \emph{all} registered members
before popping any message. Hence disjoint addressing and
unchanged admitted scope do not imply disjoint progress. A
registered observer which cannot provide a pose makes the
otherwise qualified diagnostic remain queued. This imposes an
ambient premise that must be proved at the effect cut, not
secretly carried inside the scoped evidence. In this serial replay
a fresh readiness check prevents queuing under the wrong premise.
When the test operator is allowed to remove that provisional
registration, the prepared diagnostic can complete without being
reissued. The test assigns no general right to remove a member;
at an unrepairable broker the local readiness proof cannot make
the send happen.

\paragraph{Discriminating tests.}
Execute the pinned broker and radio accounting. Compare
one held preparation against a global-generation validation and
against the direct conjunction of scoped and live requirements.
The 128 fixtures enumerate two hidden rebindings, possible source
and receiver replacement, law change, missing pose, extra receiver
and an outliving receive debit. The whole-epoch validator rejects
the original preparation in all 128; a scoped comparison plus the
live check accepts the same preparation in the two fully qualified
cases at the selected cut, and can continue it in the two readiness
counterexamples after authorized cleanup. The direct conjunction
agrees exactly. Refreshing a global token was not assigned a
latency, so the measured benefit is retention of a still-valid
preparation, not a claimed gain over all refresh strategies.
Actual selected native stalls are checked before popping; no
inference is made from the lack of a callback alone.

\paragraph{Boundary.}
This is the distinction that makes composition helpful here:
\emph{the evidence crosses, but the effect may not}. The underlying
lifting is a local, source-scoped comparison under an independently
declared live service context. The adapter is serial; physical subject
qualification, immutable named content, authority versioning,
deterministic link conditions, and stable pose availability
through dispatch are modeled. There is no transport of proofs across
unknown asynchronous process or hardware effects, no claim about an
external check-and-dispatch hook on stock SubT, and no independent
mission result.
\end{document}
